\subsection{A sharp lower bound between opposite dodecahedral vertices}
\label{sec:dodlower}

Write $\phi=(1+\sqrt5)/2$ and normalize edges to length one.
Take the twenty vertices obtained by multiplying the following vectors by
$\phi/2$. Vertices $0,\ldots,7$ are the lexicographically ordered sign
triples $(\pm1,\pm1,\pm1)$, with $-1$ first. Vertices $8,\ldots,19$ are
three blocks of four. In block $s=0,1,2$, cyclically shift
$(0,i/\phi,j\phi)$ right by $s$ places, taking the four sign choices
$(i,j)$ in lexicographic order.
The endpoints will be $V=0$ and $W=7=-V$, where the last identification
is in spatial coordinates. In particular $V=-(\phi/2)(1,1,1)$.
Here are the outward cyclic orders of the pentagons. Cyclically changing
the starting vertex of a row has no effect on the strip construction.
\[
\begin{array}{c|l@{\qquad}c|l}
0&0,16,1,14,12 &6&17,19,6,10,2\\
1&0,12,2,10,8 &7&3,11,7,19,17\\
2&0,8,4,18,16 &8&4,13,15,5,18\\
3&9,11,3,14,1 &9&8,10,6,13,4\\
4&16,18,5,9,1 &10&5,15,7,11,9\\
5&2,12,14,3,17 &11&6,19,7,15,13
\end{array}
\]
These rows can also be recovered by joining vertex pairs at unit distance
and following the five-cycles bounding their supporting planes.

Consider any finite facewise polygonal path between $V$ and $W$.
If a face occurs twice in its face list, replace the intervening portion,
including the two portions in that face, by a chord of the face. Convexity
keeps the chord on the surface and the triangle inequality bounds its
length by the former length. Repetition removal terminates since the list
gets shorter. In the resulting list, start in the last face incident to
$V$ by replacing the whole prefix ending at the exit of that face by a
chord from $V$. End at the first subsequent face incident to $W$, again
by a chord. The two kinds of faces are disjoint. Consecutive faces meet
in an edge, even when the original transition was at a vertex, since
three pentagons meet at each vertex. The threefold rotation about
$VW$ permits us to call the initial face $0$. We must thus check simple
strips starting at $0$, avoiding faces $1,2$, and stopping as soon as
one of $7,10,11$ is reached.

For completeness, the following finite computation gives a lower bound for
\emph{every} resulting strip, including any in which the developed segment
exits the strip. Put $z=\exp(2\pi i/5)$ as a planar description of an
algebraic constant. Computations use $z^4=-1-z-z^2-z^3$ and $\bar z=z^4$.
In an outward face order, the unit pentagon template has vertices
$P_k=\sum_{r=0}^{k-1}z^r$, $0\leq k\leq4$. Start with this chart on face
zero. When crossing from the known chart $T$ into a fresh face with
template positions $P(a),P(b)$ for its two shared vertex labels $a,b$,
use
\[
 T'(k)=T(a)+(T(b)-T(a))\,\overline{P(b)-P(a)}\,(P(k)-P(a)).
\]
The multiplier has unit norm, and the cyclic orders on a common edge are
opposite. Hence $T'$ is the development on the other side of that edge.
At a terminal face, record $|T'(W)-T_0(V)|^2$. The branch rule is to append
each unused face sharing two vertex labels with the current face, except
for faces $1,2$. The following table gives the pair $(a,b)$ representing
each terminal value $a+b\phi$, followed by its multiplicity.
\[
\begin{array}{rr|r@{\qquad\qquad}rr|r}
a&b&\#&a&b&\#\\ \hline
5&7&4&11&10&2\\
5&8&6&11&16&2\\
7&6&2&13&17&4\\
7&8&2&13&21&4\\
7&9&4&17&26&6\\
7&11&4&19&20&2\\
9&14&6&19&30&2
\end{array}
\]
Here is a short way to reproduce or formalize every entry. Store a
ring element as the coefficients of $1,z,z^2,z^3$. Multiply by convolution
and reduce coefficients of degrees $6,5,4$ in that order. Conjugation
sends $(a,b,c,d)$ to $(a-b,-b,d-b,c-b)$. A squared norm has the form
$(u,0,v,v)$, which is $u-v\phi$. Depth-first enumeration through the listed
faces has depth at most twelve, so the displayed recurrence with fuel
twelve is primitive recursive. The companion
\texttt{dodeca\_lower.py} implements these same rules using integers.
Every entry is at least $5+7\phi$ (the only entry with a smaller second
coefficient is $7+6\phi$, whose excess is $2-\phi>0$). A path within each
strip develops to a planar path, possibly with bends. Its length is at
least the displacement of its endpoints. This proves the required lower
bound for all finite facewise paths, and the finite approximation of
continuous paths gives the same lower bound for the continuous metric.

The minimum can be seen directly on the strip $(0,3,10)$. Its initial
endpoint is the origin and its final endpoint is
$Q=2+2z+z^2-z^3$. On the first hinge, ordered $(1,14)$, and the second,
ordered $(9,11)$, the point $tQ$ has edge parameter $t$ for
$t_1=2-\phi$ and $t_2=(2+\phi)/5$, respectively. Substitution in the
chart gives the two identities, and $0<t_1<t_2<1$. Convexity puts the
three intervening segment portions in their prescribed faces. Thus the
distance is
\[
 d(V,W)=\sqrt{5+7\phi}=\sqrt{(17+7\sqrt5)/2}.
\]
Every other vertex lies within four edges of $V$. Indeed the successive
edge-distance shells about $0$, omitting $0$ itself, are
$\{8,12,16\}$, $\{1,2,4,10,14,18\}$, $\{3,5,6,9,13,17\}$,
$\{11,15,19\}$, $\{7\}$. Since $5+7\phi>16$, any other vertex pair
except an antipodal pair is strictly closer.
